\uuid{Yejs}
\exo7id{1595}
\titre{exo7 1595}
\auteur{gineste}
\organisation{exo7}
\datecreate{2001-11-01}
\isIndication{false}
\isCorrection{false}
\chapitre{Réduction d'endomorphisme, polynôme annulateur}
\sousChapitre{Polynôme annulateur}
\cpge{mpsi-mp2i.calcul-matriciel-systemes.e, pcsi.determinants.c, ptsi.determinants.a, mpsi-mp2i.determinants.d, pcsi.calcul-matriciel-systemes-lineaires.d, ptsi.calcul-matriciel-systemes-lineaires.d, bcpst1.matrices, mp-mpi.reduction-endomorphismes.c, pc.reduction.b, psi.reduction-endomorphismes.b, pt.reduction-endomorphismes.b, mp-mpi.reduction-endomorphismes.j}
\module{Algèbre}
\niveau{L2}
\difficulte{}

\contenu{
\question{
On considére la matrice
$$ A=\left[ \begin{array}{rrr}3&-2&-1\\
2&-1&1\\
6&3&-2
 \end{array} \right].$$
Calculer son polynôme caractéristique, calculer $A^2$ et déduire de ces calculs 
et du théorème de Cayley-Hamilton l'inverse de $A.$
}
}
